% BHU PYQ -- Manually transcribed & error-corrected
\documentclass[11pt,a4paper]{article}

\usepackage[a4paper,top=2.5cm,bottom=2.5cm,left=2.8cm,right=2.8cm,headheight=14pt]{geometry}
\usepackage{amsmath,amssymb}
\usepackage[T1]{fontenc}
\usepackage[utf8]{inputenc}
\usepackage{lmodern}
\usepackage{microtype}
\usepackage{fancyhdr}
\usepackage{enumitem}
\usepackage{hyperref}
\usepackage{lastpage}
\usepackage{parskip}

\hypersetup{colorlinks=true,urlcolor=black,linkcolor=black,
  pdftitle={CHB-101: Structure and Bonding & Organic Chemistry-I -- BHU PYQ},pdfauthor={Banaras Hindu University}}

\pagestyle{fancy}\fancyhf{}
\renewcommand{\headrulewidth}{0.4pt}\renewcommand{\footrulewidth}{0.4pt}
\fancyhead[L]{\small\textit{Banaras Hindu University}}
\fancyhead[R]{\small\textit{CHB-101: Structure & Bonding & Organic Chemistry-I}}
\fancyfoot[C]{\small Page \thepage\ of \pageref{LastPage}}

\newlist{parts}{enumerate}{2}
\setlist[parts,1]{label=(\alph*),leftmargin=*,itemsep=5pt,topsep=3pt}
\setlist[parts,2]{label=(\roman*),leftmargin=*,itemsep=3pt,topsep=2pt}
\newcommand{\pts}[1]{\hfill\textit{[#1]}}

\begin{document}

\begin{center}
  {\large\bfseries BANARAS HINDU UNIVERSITY}\\[2pt]
  {\normalsize B.Sc. (Hons.) Semester I Examination, 2016-17}\\[6pt]
  \rule{0.8\linewidth}{0.5pt}\\[6pt]
  {\large\bfseries Paper No. CHB-101: Structure \& Bonding and Organic Chemistry-I}\\[4pt]
  {\normalsize Time: 3 Hours\hfill Maximum Marks: 70}
\end{center}
\rule{\linewidth}{0.4pt}
\small
\noindent\textit{Instructions: Answer five questions in total, including Question~1 from each section which is compulsory. Use separate answer books for Section~A and Section~B.}
\normalsize
\bigskip

\begin{center}
  {\large\bfseries SECTION A}\\[2pt]
  {\normalsize (Structure and Bonding -- Marks: 35)}
\end{center}
\rule{0.4\linewidth}{0.2pt}
\smallskip

\noindent\textbf{Question 1.} Answer all parts of the following: \pts{5 $\times$ 2 = 10}
\begin{parts}
  \item Depict the radial distribution curve for a $4\text{d}$ orbital and mention the number of radial nodes in it.
  \item Why are Group 1 metals quite soft as compared to transition metals? Explain.
  \item What is the radius ratio rule? In what way does it help in predicting the coordination number of ionic crystals?
  \item Arrange the following in order of increasing covalent character, giving appropriate reason: $\text{NaCl}$, $\text{MgCl}_2$, $\text{AlCl}_3$.
  \item Calculate the formal charge of each atom in the nitrite ion ($\text{NO}_2^-$).
\end{parts}

\medskip\rule{\linewidth}{0.2pt}

\noindent\textbf{Question 2.}
\begin{parts}
  \item State Schr\"odinger wave equation and discuss the physical significance of $\psi$ and $\psi^2$. \pts{6}
  \item Draw the molecular orbital energy level diagram for $\text{NO}$ and mention the HOMO and LUMO. Calculate its bond order. \pts{6}
\end{parts}

\medskip\rule{\linewidth}{0.2pt}

\noindent\textbf{Question 3.}
\begin{parts}
  \item With the help of suitable examples, discuss van der Waals interactions. Are these stronger or weaker than hydrogen bonding? \pts{6}
  \item Discuss Slater's rules for calculating the effective nuclear charge ($Z_{\text{eff}}$) and calculate $Z_{\text{eff}}$ for a $3\text{d}$ electron in a zinc atom ($Z = 30$). \pts{6}
\end{parts}

\medskip\rule{\linewidth}{0.2pt}

\noindent\textbf{Question 4.}
\begin{parts}
  \item Explain why beryllium exists in the metallic form and not as a diatomic gas molecule ($\text{Be}_2$) on the basis of Molecular Orbital Theory. \pts{6}
  \item Using VSEPR theory, predict the geometries and draw the structures of the following: (i) $\text{SF}_4$, (ii) $\text{ClF}_3$, and (iii) $\text{XeF}_2$. \pts{6}
\end{parts}

\newpage
\begin{center}
  {\large\bfseries SECTION B}\\[2pt]
  {\normalsize (Organic Chemistry-I -- Marks: 35)}
\end{center}
\rule{0.4\linewidth}{0.2pt}
\smallskip

\noindent\textbf{Question 5.} Answer all parts of the following: \pts{5 $\times$ 2 = 10}
\begin{parts}
  \item Why is the enol form of ethyl acetoacetate more stable than that of simple acetone?
  \item Explain why methyl radical ($\cdot\text{CH}_3$) has a planar geometry whereas trifluoromethyl radical ($\cdot\text{CF}_3$) has a pyramidal geometry.
  \item Draw the chair and boat conformations of cyclohexane and explain their relative stability.
  \item What are enantiomers and diastereomers? Explain with one suitable example of each.
  \item Briefly discuss the mechanism of $\text{S}_{\text{N}}2$ reaction with a suitable example.
\end{parts}

\medskip\rule{\linewidth}{0.2pt}

\noindent\textbf{Question 6.}
\begin{parts}
  \item How will you prepare the following compounds starting from ethyl acetoacetate? Write stepwise mechanisms: \pts{3 $\times$ 4 = 12}
  \begin{parts}
    \item Ethyl methyl ketone
    \item $n$-Valeric acid
    \item Succinic acid
    \item Cyclopropanecarboxylic acid
  \end{parts}
\end{parts}
\hfill \textbf{OR}
\begin{parts}
  \item How are the following compounds prepared from Grignard reagent ($\text{RMgBr}$)? Write reactions with suitable mechanisms: \pts{3 $\times$ 4 = 12}
  \begin{parts}
    \item Primary alcohol
    \item Carboxylic acid
    \item Ketone
    \item Alkane
  \end{parts}
\end{parts}

\medskip\rule{\linewidth}{0.2pt}

\noindent\textbf{Question 7.} Explain any four of the following: \pts{4 $\times$ 3 = 12}
\begin{parts}
  \item $1,3$-Butadiene gives $1,2$- as well as $1,4$-addition products with bromine.
  \item Dehydration of alcohols generally needs an acid catalyst.
  \item \textit{cis}-Hydroxylation of $2$-butene with cold alkaline $\text{KMnO}_4$ gives \textit{meso}-$2,3$-butanediol.
  \item Chlorination of methane in the presence of sunlight is a free-radical chain reaction.
  \item $\text{S}_{\text{N}}1$ reactions are accompanied by racemization at the chiral center.
\end{parts}

\medskip\rule{\linewidth}{0.2pt}

\noindent\textbf{Question 8.}
\begin{parts}
  \item Draw structural formulae for each of the following pairs of compounds: \pts{4}
  \begin{parts}
    \item $(E)$- and $(Z)$-$2$-pentene
    \item $(E)$- and $(Z)$-$2$-phenyl-$2$-butene
  \end{parts}
  \item How will you carry out the following conversions? Write mechanisms in each case: \pts{4 $\times$ 2 = 8}
  \begin{parts}
    \item Cyclohexene to \textit{trans}-$1,2$-cyclohexanediol
    \item Propene to $1,2$-dibromopropane
  \end{parts}
\end{parts}

\vfill
\rule{\linewidth}{0.4pt}
\begin{center}\small
  \href{https://bhu-exam.vercel.app/}{Click here to visit our website}\\[4pt]
  \href{https://me-aryan.vercel.app/}{Click here to visit our contributor}\\[4pt]
  \href{https://quashan-me.vercel.app/}{Click here to visit our contributor}
\end{center}

\end{document}
